\section{Coherence, representation and conditioning of lower previsions}
\label{sec:annex-p15-imprecise-probability}

\begin{annexitem}{annex:p15.ip.stakes}{thm:p15-ip-stakes}{Passage from integer portfolios to real stakes}
\begin{proof}
Integer stakes are special real stakes. Conversely rational nonnegative
stakes can be cleared of denominators; repetitions of a gamble then
produce the stated integer portfolio. For fixed finitely many gambles,
the portfolio is continuous in its coefficients in the supremum norm:
\[
 \left\|\sum_i(s_i-t_i)G(f_i)\right\|_\infty
       \leq\sum_i|s_i-t_i|\,\|G(f_i)\|_\infty.
\]
The supremum of a bounded function changes by at most its supremum-norm
distance. Approximate any finite real-stake portfolio by rational ones,
including the single selling coefficient in the coherent case.
Its supremum is the limit of their nonnegative suprema, hence is
nonnegative. The empty portfolio is zero.
\end{proof}
\end{annexitem}

\begin{annexitem}{annex:p15.ip.natural-extension}{thm:p15-ip-natural-extension}{Finite extension and the exact original-domain criterion}
\begin{proof}
Write \(S=\sum_i s_iG(f_i)\) for a buying portfolio. The empty
portfolio makes every \(\alpha\leq\inf g\) feasible. If a feasible
\(\alpha>\sup g\) existed, \(S\leq g-\alpha\) would have negative
supremum, violating avoidance of sure loss. This proves finiteness
and the bounds. It also proves \(\underline E(c)=c\).

Pointwise order gives monotonicity. Translating a feasible bound for
\(g\) by \(c\) gives exactly a feasible bound for \(g+c\), proving
cash translation in both directions. For \(t>0\), multiplication and
division of both the gamble and all portfolio coefficients prove
positive homogeneity; \(t=0\) uses \(\underline E(0)=0\).
Adding feasible inequalities for \(g,h\) gives a feasible inequality
for \(g+h\). Choose bounds within any positive tolerance of the two
suprema and let the tolerance decrease to zero, proving superadditivity.
The pointwise inequalities
\(g\leq h+\|g-h\|_\infty\) and its reverse give the Lipschitz bound.

We prove full-domain coherence without presupposing representation.
Each centered gain \(g-\underline E(g)\) has lower prevision zero
by cash translation. A positive buying portfolio \(H\) of these
gains has \(\underline E(H)\geq0\) by homogeneity and superadditivity.
If \(H-s_0[g_0-\underline E(g_0)]\) had supremum \(-\delta<0\),
then monotonicity and cash translation would imply
\[
 0\leq\underline E(H)
 \leq\underline E(s_0[g_0-\underline E(g_0)]-\delta)
 =-\delta,
\]
a contradiction. This includes \(s_0=0\) and proves avoidance and
the one-selling-gamble condition.

Taking a one-unit portfolio of \(f\in\mathcal K\) gives
\(\underline E(f)\geq\underline P(f)\). If the input is coherent
and \(\underline E(f_0)>\underline P(f_0)\), choose a feasible
\(\alpha>\underline P(f_0)\). The defining inequality then gives
\[
 \sup[S-(f_0-\underline P(f_0))]
                 \leq\underline P(f_0)-\alpha<0,
\]
contrary to coherence. Thus the extension agrees with every assessment.
Conversely restriction of a coherent full-domain functional preserves
all the required portfolio inequalities, so agreement implies coherence.

Let \(J\) be any coherent full-domain lower prevision dominating the
input. For a feasible inequality, each \(J(f_i-\underline P(f_i))\)
is nonnegative. Consequently \(J(g-\alpha)\geq0\), and cash
translation gives \(J(g)\geq\alpha\). Taking the supremum gives
\(J(g)\geq\underline E(g)\). The functional properties for a
coherent \(J\) follow independently in the next item, so this use
does not assume a representation theorem.

Finally if \(S\) is an initial portfolio with \(\sup S=-\delta<0\),
then \(g-\alpha\geq tS\) whenever
\(\alpha\leq\inf g+t\delta\). Arbitrarily large \(t\) make the
natural extension infinite on every \(g\). A coherent finite
dominating functional would give \(J(S)\geq0\) and
\(J(S)\leq-\delta\), so cannot exist.
\end{proof}
\end{annexitem}

\begin{annexitem}{annex:p15.ip.linear-domain}{thm:p15-ip-linear-domain}{The vector-domain characterization}
\begin{proof}
The function-space identifications used in the domain definition follow
by quantizing a bounded function with an arbitrarily fine finite grid.
On the power-set algebra every grid cell is an available event; on a
sigma algebra the cells belong to it whenever the function is measurable.
The resulting simple functions converge uniformly. Conversely a uniform
limit \(g\) of measurable simple \(g_n\) is measurable, since
\[
 \{g>r\}=\bigcup_{k\geq1}\bigcup_{N\geq1}
                 \bigcap_{n\geq N}\{g_n>r+1/k\}.
\]
For the forward inclusion choose \(1/k\) below half the positive
gap and use convergence; the reverse inclusion follows by taking limits.
Thus the asserted full-domain specializations preserve measurability.

Suppose first that the portfolio criterion holds on \(B\).
For any \(g\), the selling-only portfolio implies \(I(g)\geq\inf g\).
The buying-only portfolio gives \(I(g)\leq\sup g\); hence constants
are fixed. The two admissible centered portfolios
\(t[g-I(g)]-[tg-I(tg)]\) and
\([tg-I(tg)]-t[g-I(g)]\) are opposite constants.
If \(I(tg)\ne tI(g)\), one has a strictly negative gain;
therefore equality holds for \(t>0\). If
\(I(g+h)<I(g)+I(h)\), buying \(g,h\) and selling their sum also
gives a strictly negative constant gain. These prove the three properties.

Conversely assume the three displayed properties. Homogeneity at
\(g=0,t=2\) gives \(I(0)=0\). Superadditivity applied to \(g,-g\)
gives \(I(g)+I(-g)\leq0\), and the lower bounds yield the corresponding
upper bounds. Constants are therefore fixed. Monotonicity follows from
\(I(h)\geq I(g)+I(h-g)\geq I(g)\) when \(h\geq g\).
Superadditivity with a constant in both directions gives
\(I(g+c)=I(g)+c\). Homogeneity extends to \(t=0\), and these
properties also give the Lipschitz estimate exactly as in the preceding
item. Its centered-portfolio argument proves coherence, using only
the properties just established. This also justifies the use of those
properties for a coherent \(J\) in the least-extension argument.
\end{proof}
\end{annexitem}

\begin{annexitem}{annex:p15.ip.representation}{thm:p15-ip-representation}{Compatible charges on the full function space}
\begin{proof}
First suppose the input avoids sure loss, so that its finite superlinear
natural extension has been constructed. Apply the general dominated
extension theorem \cref{thm:p17-mmeu-hahn-banach} to
\(p=-\underline E\). For a fixed nonzero \(g\), start on its linear
span with \(\ell(tg)=-t\underline E(g)\). For \(t\geq0\)
this equals \(p(tg)\); for \(t<0\), domination follows from
\(\underline E(g)+\underline E(-g)\leq0\). For \(g=0\)
start on the zero subspace. The extended linear \(\ell\leq p\)
gives \(L=-\ell\geq\underline E\), with
\(L(g)=\underline E(g)\). Lower bounds on positive gambles make
\(L\) positive. The two constants give \(L(1)=1\), and positivity
implies \(|L(f)|\leq\|f\|_\infty\). Thus it is a precise model
and belongs to \(\mathcal M(\underline P)\).

Conversely any compatible \(L\) assigns a buying portfolio \(S\)
expectation \(L(S)\geq0\), while \(L(S)\leq\sup S\).
This proves avoidance of sure loss. For every defining feasible bound,
\(L(g)-\alpha\geq L(S)\geq0\), so every compatible \(L\)
dominates \(\underline E\). Conversely domination of
\(\underline E\) implies domination of the input. This proves
equality of the two compatible sets. The supporting \(L\) at each
\(g\) proves attainment and the lower-envelope formula. Agreement
on the assessed domain is exactly the previously proved coherence
criterion.

The precise models correspond to probabilities on the stated algebra
by \cref{thm:p17-mmeu-charges-compactness}. Each inequality
\(L(f)\geq\underline P(f)\) defines a closed subset of that compact
space because bounded-gamble integration is continuous. Their intersection
is closed, nonempty and convex, hence compact.

For full-domain uniqueness, let nonempty closed convex \(C\) represent
\(I\) on every member of \(B\). Certainly \(C\subseteq\mathcal M(I)\).
If \(L_0\in\mathcal M(I)\setminus C\), the finite-coordinate
separation theorem \cref{thm:p17-mmeu-finite-separation} supplies a
simple \(g\) with \(L_0(g)<\min_{L\in C}L(g)=I(g)\), a
contradiction. Thus \(C=\mathcal M(I)\). An arbitrary partial
assessment domain need not contain this separating gamble, which is
why the same uniqueness is not asserted there.
\end{proof}
\end{annexitem}

\begin{annexitem}{annex:p15.ip.dual}{thm:p15-ip-dual}{Conjugacy and transitive extension}
\begin{proof}
Negating the minimum of \(L(-g)\) gives the maximum of \(L(g)\),
with attainment from the same supporting functional argument. Minimum
is at most maximum. Cash translation and
\(-I_A=I_{A^c}-1\) yield
\(-\underline E(-I_A)=1-\underline E(I_{A^c})\).

Let \(\mathcal K\subseteq\mathcal K_1\subseteq\mathcal K_2\subseteq B\),
and restrict the original natural extension to \(\mathcal K_1\).
Every compatible original \(L\) dominates this restriction, by the
proved extension formula. Conversely every model dominating that
restriction dominates the original assessments. The compatible sets
are identical. Their lower envelopes on \(\mathcal K_2\) are
therefore identical, proving transitivity. The empty initial domain
gives the vacuous functional and satisfies the same argument.
\end{proof}
\end{annexitem}

\begin{annexitem}{annex:p15.ip.countable-additivity}{thm:p15-ip-countable-additivity}{Upper continuity and the finitely additive counterexample}
\begin{proof}
A finitely additive probability on a sigma algebra is countably additive
exactly when \(L(A_n)\to0\) for every \(A_n\downarrow\varnothing\).
For the nontrivial implication, given disjoint \(E_i\) let
\(E=\bigcup_iE_i\). The remainders
\(A_n=E\setminus\bigcup_{i\leq n}E_i\) decrease to the empty set.
Finite additivity and the limit give
\(L(E)=\sum_iL(E_i)\). The reverse implication follows by writing
the decreasing first set as its disjoint successive differences.

If upper continuity holds, each compatible \(L\) has
\(0\leq L(A_n)\leq\overline E(A_n)\to0\), proving its
countable additivity. Conversely suppose all compatible models are
countably additive. For \(\epsilon>0\), the open sets
\(U_n=\{L:L(A_n)<\epsilon\}\) increase and cover the compact
compatible set. A finite subcover, followed by its largest index,
gives \(L(A_N)<\epsilon\) for every \(L\). The maximum is
attained, so \(\overline E(A_N)<\epsilon\); monotonicity then
proves the stated limit.

On \(\mathbb N\), extend the cofinite filter to an ultrafilter
using the supplied Choice principle. It contains no finite set.
The probability \(L_{\mathcal U}(A)=1\) for \(A\in\mathcal U\)
and zero otherwise is finitely additive: a disjoint union is in the
ultrafilter exactly when one of its two members is. It gives every
singleton zero and the whole space one, contradicting countable
additivity. Its bounded-gamble integral is a coherent linear prevision
by the representation already supplied in
\cref{thm:p17-mmeu-domain-boundaries}.

The vacuous lower prevision \(\inf g\) is coherent by the vector-domain
criterion. Every precise model dominates it, including this ultrafilter
model. If \(A_n\downarrow\varnothing\), any chosen point eventually
lies outside \(A_n\), so \(\inf I_{A_n}=0\) eventually.
Thus lower continuity holds without making all compatible models
countably additive.
\end{proof}
\end{annexitem}

\begin{annexitem}{annex:p15.ip.event-separation}{thm:p15-ip-event-separation}{All event bounds and the separating utility gamble}
\begin{proof}
For both vertex lists every coordinate attains minimum \(1/10\)
and maximum \(3/5\). Removing \((1,3,6)/10\) removes neither
extremum: coordinate 1 still attains its minimum at \((1,6,3)/10\),
and coordinate 3 still attains its maximum at \((3,1,6)/10\).
The other extrema also remain in the list. Linear evaluation on a
convex hull has the same extrema as evaluation on its generating
vertices. This gives all singleton lower and upper probabilities.
For a pair use the complement singleton: its lower probability is
\(1-3/5=2/5\), and its upper probability is \(1-1/10=9/10\).
Empty and full events give zero and one. This exhausts all eight events.

In permutation order
\((1,3,6),(1,6,3),(3,1,6),(3,6,1),(6,1,3),(6,3,1)\),
each divided by ten, the expectations of \((2,1,0)\) are
\[
                  1/2,\quad4/5,\quad7/10,\quad6/5,\quad13/10,\quad3/2.
\]
The full hull has minimum \(1/2\); the reduced hull has minimum
\(7/10\). Both full gamble envelopes are coherent, as minima of
linear expectations. The certain payoff has expectation \(3/5\)
under every model. A maxmin comparison therefore chooses the certain
payoff for \(C\) and the gamble for \(D\). The differences from
the constant also have lower expectations \(-1/10\) and \(1/10\),
respectively. The conclusion is under the stated decision rules, not
a completeness requirement imposed by coherence itself.
\end{proof}
\end{annexitem}

\begin{annexitem}{annex:p15.ip.possibility-bridge}{thm:p15-ip-possibility-bridge}{The nested mass calculation}
\begin{proof}
The ordered outcomes exist because the space is finite; break equal-score
ties arbitrarily. Each difference is nonnegative and their telescoping
sum is \(\pi_1-\pi_{n+1}=1\); no empty focal set receives mass.
If \(A\ne\Omega\), let \(j\) be its first omitted outcome in this
order. Exactly the sets \(B_i\) with \(i<j\) are contained in
\(A\), so their mass sum is
\(\pi_1-\pi_j=1-\max_{\omega\notin A}\pi(\omega)\).
For \(A=\Omega\) the sum is one, giving the same necessity value
under the empty-maximum convention. For \(A\ne\varnothing\),
let \(k\) be its first included outcome. Exactly the sets with
\(i\geq k\) intersect it, giving
\(\sum_{i=k}^n(\pi_i-\pi_{i+1})=\pi_k=
\max_{\omega\in A}\pi(\omega)\). For empty \(A\) both
the plausibility and possibility value are zero. This proves the two
event identities, including ties and zero scores.

Apply the fully proved allocated-envelope result
\cref{thm:p14-evidence-finite-measures} to these masses. Expected values
under its nonempty allocated probability family have a coherent lower
envelope by the vector-domain criterion, whose restriction to indicators
is the specified necessity measure; the conjugate upper restriction is
the possibility measure. This constructs the advertised extension and
uses no claim that all coherent lower probabilities are of this nested
form.
\end{proof}
\end{annexitem}

\begin{annexitem}{annex:p15.ip.conditional-models}{thm:p15-ip-conditional-models}{One-event conditioning with every zero case retained}
\begin{proof}
A supported precise \(Q\) satisfies \(Q(g)=Q(I_Ag)\): the
absolute value of \((1-I_A)g\) is at most
\(\|g\|_\infty(1-I_A)\), whose expectation is zero.
If \(L(A)>0\), the functional
\(Q_L(g)=L(I_Ag)/L(A)\) is positive, linear, normalized and supported
on \(A\). If \(L(A)=0\), positivity bounds
\(|L(I_Ag)|\leq\|g\|_\infty L(A)=0\), so every supported \(Q\)
is Bayes-compatible. There is always such a \(Q\), for example a
point mass at any point of nonempty \(A\).

Every nonempty collection of supported precise models has a coherent
lower envelope on the restricted gamble space: the infimum is bounded
below by \(\inf_A g\), positive homogeneous, superadditive and fixes
constants. The vector-domain criterion applies. This proves separate
coherence of both updates, including the stated vacuous convention.

If \(\underline E(A)>0\), every original model is positive on \(A\)
and has its unique Bayesian conditional. All pairs therefore give the
same natural and regular family. Evaluation of \(Q_L(g)\) is a
continuous ratio of bounded-integral coordinates with denominator bounded
away from zero. Compactness of \(\mathcal M\) gives attainment.
If \(\underline E(A)=0\), compactness supplies a model \(L_0\)
with \(L_0(A)=0\). Every supported \(Q\) occurs in a compatible
pair with it; the conditional infimum is then \(\inf_A g\), as
point masses approach that infimum. If \(\overline E(A)>0\), at
least one model \(L_1\) has positive \(A\)-mass, so regular
conditioning has a nonempty Bayesian family. If the upper probability
is zero, no such family exists and the declared vacuous value applies.

We justify joint coherence on the stated single-event transaction
domain. A precise Bayes-compatible pair has no uniformly losing
portfolio. For a portfolio involving unconditional transactions,
\(L\) assigns every centered gain expectation zero, including
conditional gains because of the Bayes identity. Thus its supremum
on \(\Omega\) is nonnegative. For a portfolio of conditional
transactions only, \(Q\) assigns expectation zero and is supported
on \(A\), so its supremum on \(A\) is nonnegative.

For lower-envelope buying gains these expectations are nonnegative.
For the single sold gain, choose a pair whose corresponding precise
price is within \(\epsilon\) of its lower-envelope price. In the
unconditional-sale case use \(L\); in the conditional-sale case
use its \(Q\). The portfolio expectation is at least
\(-s_0\epsilon\) on the appropriate union. For a conditional-only
portfolio use \(Q\); for a portfolio also involving unconditional
transactions use \(L\), where the conditional error is multiplied
by \(L(A)\leq1\). Let \(\epsilon\) decrease to zero. This proves
all the finite lower-envelope portfolio inequalities, along with the
required separate conditional tests.

For natural conditioning the family projects onto all original \(L\),
so its unconditional envelope is unchanged. For regular conditioning
when \(\overline E(A)>0\), every zero-\(A\) model is approached by
\((1-t)L_0+tL_1\), \(t\downarrow0\), inside the convex compatible
set. Its positive-\(A\) part is therefore dense in that set and has
the same unconditional infimum for every continuous bounded-gamble
evaluation. The same portfolio proof consequently uses the unchanged
unconditional lower prevision. This argument proves a single-event
claim, not an envelope theorem for arbitrary infinite partitions.
\end{proof}
\end{annexitem}

\begin{annexitem}{annex:p15.ip.generalized-bayes}{thm:p15-ip-generalized-bayes}{A unique root only at positive lower probability}
\begin{proof}
Put \(F(q)=\underline E(I_A[g-q])\). For a positive-\(A\) model,
\[
 L(I_A[g-q])=L(A)[Q_L(g)-q].
\]
Suppose first \(a=\underline E(A)>0\). The conditional minimum
\(r\) is attained. At \(q=r\), all these products are nonnegative
and at least one is zero, so \(F(r)=0\). For \(\delta>0\),
\[
 F(q)-\overline E(A)\delta
       \leq F(q+\delta)\leq F(q)-a\delta.
\]
These inequalities follow by bounding each \(L(A)\) before taking
infima; hence the function is strictly decreasing and continuous.
Its sole root is \(r\).

If \(a=0<\overline E(A)\), a zero-\(A\) model evaluates every
argument of \(F\) as zero. For \(q\leq r\), all positive-model
products are nonnegative, so \(F(q)=0\). For \(q>r\), the
definition of the conditional infimum supplies a positive model with
\(Q_L(g)<q\); its product is negative and thus \(F(q)<0\).
The root set is exactly \(( -\infty,r]\). In particular a root
equation alone cannot distinguish the natural vacuous conditional
from a stronger selected value at lower probability zero.
\end{proof}
\end{annexitem}

\begin{annexitem}{annex:p15.ip.conditional-bounds}{thm:p15-ip-conditional-bounds}{Two portfolios force the generalized-Bayes identity}
\begin{proof}
Fix \(g\) and put \(q=C(g\mid A)\), \(h=I_A(g-q)\).
A one-unit conditional buying gain is \(h\); the unconditional
selling gain for the assessed full-domain gamble \(h\) is
\(-(h-\underline E(h))\). Their sum is the constant
\(\underline E(h)\). Since an unconditional transaction is present,
joint coherence takes its supremum on \(\Omega\), giving
\(\underline E(h)\geq0\). Conversely buy \(h\) unconditionally
and sell \(g\) conditionally. The net gain is
\((h-\underline E(h))-h=-\underline E(h)\), so coherence gives
\(\underline E(h)\leq0\). This proves the identity without
assuming a set-of-probabilities interpretation of the subject's prices.

At positive lower probability the unique-root theorem forces
\(q=\underline R(g\mid A)\). At zero lower and positive upper
probability, its full root set gives the upper bound on \(q\).
Separate coherence on the restricted gamble space gives
\(q\geq\inf_A g\). The two limiting functionals themselves are
jointly coherent by \cref{thm:p15-ip-conditional-models}. Their
pointwise extremal character therefore follows from the bounds just
proved. No conclusion about arbitrary infinite partitions or about
additional accepted transactions at exact price endpoints was used.
\end{proof}
\end{annexitem}

\begin{annexitem}{annex:p15.ip.independence-boundary}{thm:p15-ip-independence-boundary}{The full strong product and its two conditional envelopes}
\begin{proof}
For each generating product with \(p=p_X(1)\in[2/5,3/5]\)
and \(q=p_Y(1)\in[0,1]\), the masses of \(A\) and
\(\{X=1\}\cap A\) are \(q\) and \(pq\).
For a convex mixture indexed by \(j\), a positive-\(A\) conditional
success probability is therefore
\[
           \frac{\sum_j w_jp_jq_j}{\sum_j w_jq_j},
\]
a convex combination of \(p_j\) with weights
\(w_jq_j/\sum_k w_kq_k\). It belongs to \([2/5,3/5]\).
The same inequalities \((2/5)L(A)\leq L(\{X=1\}\cap A)
\leq(3/5)L(A)\) persist under closure. Products with \(p=2/5\)
and \(q>0\) attain the regular lower endpoint. Marginal success
probabilities are mixtures of the \(p_j\), with the same attainable
endpoints. Products with \(q=0\) and \(q=1\) show that the
lower and upper probability of \(A\) are zero and one.

The zero-\(A\) product leaves every supported conditional model
Bayes-compatible. In particular the conditional point mass at
\((X,Y)=(0,1)\) gives conditional lower success probability zero;
nonnegativity supplies the matching lower bound. This proves natural
vacuity for the success indicator and the stated distinction from
regular conditioning. The construction uses the full declared product
outcome space and specifies the independence model; no independence
assumption is inferred from an event interval alone.
\end{proof}
\end{annexitem}
